\subsection{Implementierung der Medikamentenverwaltung (5123060)}
Auch die Medikamentenverwaltung bildet eine zentrale fachliche Komponente innerhalb des Hospital Management Systems.
Sie umfasst die Verwaltung von Diagnosen (\textit{Diagnoses}), sowie darauf aufbauende konkrete Medikationsverläufe (\textit{Medications}), Medikamentenstammdaten (\textit{Drugs}) und Dosierungen (\textit{Doses}).
Durch diese Aufteilung wird eine klare Trennung zwischen statischen Informationen (z. B. Wirkstoffe), parametrisierten Verabreichungen und patientenspezifischen Behandlungsverläufen erreicht.

\subsubsection{API-Design und Datenübertragung}
Die REST-API folgt einer konsistenten Struktur entlang der Ressourcen:

\begin{itemize}
\item \textbf{Lesende Zugriffe (GET):}
Die API liefert bei komplexeren Entitäten (z. B. \textit{Medications} oder \textit{Diagnoses}) verschachtelte Objekte zurück, um kontextrelevante Informationen direkt verfügbar zu machen.
Dadurch werden zusätzliche Requests reduziert, was insbesondere bei der Anzeige medizinischer Verläufe von Vorteil ist.

\item \textbf{Schreibzugriffe (POST/PUT):}
Für schreibende Operationen werden bewusst schlanke Request-Objekte verwendet, die ausschließlich auf ID-Referenzen basieren (z. B. \texttt{drugId}, \texttt{doseId}). 
Dies reduziert die Payload-Größe und verhindert Inkonsistenzen durch redundante Datenübertragung.

\end{itemize}

\textbf{Beispiel für eine Response (\texttt{\detokenize{GET /api/medications/{MEDICATION_ID}}}):}

\begin{lstlisting}
{
"(*\color{keyblue}id*)": (*\color{numorange}1*),
"(*\color{keyblue}dose*)": {
"(*\color{keyblue}id*)": (*\color{numorange}324*),
"(*\color{keyblue}unit*)": "(*\color{valgreen}G*)",
"(*\color{keyblue}amount*)": (*\color{numorange}2*),
"(*\color{keyblue}frequency*)": "(*\color{valgreen}X\_WEEKLY*)",
"(*\color{keyblue}frequencyAmount*)": (*\color{numorange}7*)
},
"(*\color{keyblue}drug*)": {
"(*\color{keyblue}id*)": (*\color{numorange}635*),
"(*\color{keyblue}stock*)": (*\color{numorange}1384*),
"(*\color{keyblue}name*)": "(*\color{valgreen}Paracetamol*)",
"(*\color{keyblue}activeIngredient*)": "(*\color{valgreen}Paracetamol*)",
"(*\color{keyblue}type*)": "(*\color{valgreen}SYRUP*)"
},
"(*\color{keyblue}started*)": "(*\color{valgreen}2024-03-01*)",
"(*\color{keyblue}ended*)": null
}
\end{lstlisting}

\subsubsection{Filter- und Suchmechanismen}
Alle Ressourcen unterstützen flexible Filtermechanismen über Query-Parameter.
Diese erlauben eine präzise und performante Einschränkung der Datenmengen:

\begin{itemize}
\item \textbf{Drugs:} Filter nach Name, Namensbestandteilen, Wirkstoff und Typ.
\item \textbf{Doses:} Filter nach Einheit, Frequenz sowie Mengenangaben.
\item \textbf{Medications:} Filter nach Medikament, Typ, Dosierungseinheit sowie zeitlichen Einschränkungen (Startdatum).
\item \textbf{Diagnoses:} Erweiterte Filtermöglichkeiten inkl. Krankheit, behandelndem Arzt, Patient sowie Zeitraum.
\end{itemize}

Zur Sicherstellung konsistenter Suchlogik werden bestimmte Parameterkombinationen bewusst ausgeschlossen.
Beispielsweise ist die gleichzeitige Verwendung von exakter und unscharfer Suche (z. B. \texttt{name} und \texttt{nameContains}) nicht erlaubt und führt zu einem \texttt{400 BAD REQUEST}.

Bei Medikationen wird außerdem sichergestellt, dass das Enddatum nicht vor dem Startdatum liegt.

\subsubsection{Fachliche Besonderheiten}
Ein zentrales fachliches Konzept ist die Trennung zwischen \textbf{aktiven} und \textbf{beendeten} Medikationen.
Eine Medikation gilt als aktiv, solange kein Enddatum gesetzt ist.
Diese Logik wird direkt durch den API-Parameter \texttt{active} abgebildet und ermöglicht eine effiziente Abfrage aktuell laufender Behandlungen.

Das selbe gilt für Diagnosen, die ebenfalls über den \texttt{active}-Parameter gefiltert werden können.

Darüber hinaus wird perspektivisch eine Erweiterung der Medikamentenlogik angestrebt, bei der der Lagerbestand (\textit{stock}) in Relation zu aktiven Medikationen gesetzt wird, um kritische Bestände frühzeitig zu erkennen.

\subsubsection{Beispielhafte API-Nutzung}
Die folgenden Beispiele verdeutlichen typische Anwendungsfälle:

\begin{itemize}
    \item \textbf{Abruf von Daten:}
    \begin{itemize}
        \item \texttt{GET /api/drugs?activeIngredient=Insulin}: Filterung nach Wirkstoff.
        \item \texttt{GET /api/medications?drugType=SYRUP\&startedAfter=2024-01-01}: Abruf aller Medikationen mit einem Medikament des Typs SYRUP, die nach dem angegebenen Datum begonnen haben.
        \item \texttt{GET /api/diagnoses?diseaseContains=Diabetes}: Suche nach Diagnosen anhand eines Teilstrings.
    \end{itemize}

    \item \textbf{Erstellung neuer Einträge:}
    \begin{itemize}
        \item \texttt{POST /api/medications}: Erstellung einer neuen Medikation unter Verwendung von \texttt{doseId} und \texttt{drugId}.
    \end{itemize}

    \item \textbf{Aktualisierung bestehender Daten:}
    \begin{itemize}
        \item \texttt{PUT \detokenize{/api/medications/{MEDICATION_ID}}}: Setzen eines Enddatums zur Beendigung einer Medikation.
    \end{itemize}
\end{itemize}


Die gewählte Architektur stellt sicher, dass sowohl einfache als auch komplexe medizinische Zusammenhänge effizient abgebildet werden können, ohne die Performance oder Wartbarkeit des Systems zu beeinträchtigen.
